% !TEX program = xelatex
% ============================================================
%  تقويم — Calendar Template
%  تقويم شهري بشبكة سبعة أعمدة (السبت..الجمعة) مع مساحة للملاحظات.
%  صفحة واحدة لكل شهر، بتنسيق أفقي (Landscape).
%  Compile with: xelatex main.tex (run twice)
% ============================================================
%  Features:
%    - Landscape A4 wall-calendar layout
%    - TikZ 7-column grid, Arabic week (السبت..الجمعة)
%    - Month + year header placeholder
%    - Day cells with number and notes space
%    - Sample month included (reusable per month)
% ============================================================

\documentclass[11pt,a4paper]{article}

% ---- Geometry (landscape, tighter margins for wall calendar) ----
\usepackage[a4paper,landscape,margin=1.2cm,top=1.2cm,bottom=1.2cm,headheight=16pt]{geometry}

% ---- Core packages (order matters for bidi) ----
\usepackage{graphicx}
\usepackage{array}
\usepackage{tabularx}
\usepackage{setspace}
\usepackage{xcolor}
\usepackage{tikz}
\usetikzlibrary{calc,positioning}

% ---- RTL / Arabic language support (must load before hyperref) ----
\usepackage{bidi}
\usepackage[no-math]{fontspec}
\usepackage[quiet]{polyglossia}

% ---- Fonts ----
\setmainfont[Scale=1]{Amiri-Regular.ttf}
\newfontfamily\arabicfont[Script=Arabic,Scale=0.95]{Amiri-Regular.ttf}
\newfontfamily\englishfont[Scale=0.95]{TeX Gyre Termes}

% ---- Languages ----
\setdefaultlanguage[calendar=gregorian,hijricorrection=1]{arabic}
\setotherlanguage[variant=british]{english}

% ---- Hyperref (load after polyglossia/bidi) ----
\usepackage[breaklinks,hidelinks]{hyperref}
\hypersetup{
  pdftitle={تقويم},
  pdfauthor={الاسم}
}

% ---- Colors ----
\definecolor{calaccent}{HTML}{2F5D8A}
\definecolor{calhead}{HTML}{1F3F5F}
\definecolor{calweekend}{HTML}{B23A3A}
\definecolor{callight}{HTML}{F2F5F9}

% ---- Line spacing ----
\setstretch{1.2}

\pagestyle{empty}

% ---- Calendar grid parameters ----
% 7 columns x 6 rows. Usable width ~ 27.3cm, height ~ 18.6cm.
\newcommand{\cellw}{3.8}   % column width (cm)
\newcommand{\cellh}{2.3}   % row height (cm)
\newcommand{\gridw}{26.6}  % 7 * 3.8
\newcommand{\gridh}{13.8}  % 6 * 2.3

% Day headers (Arabic week starts Saturday). RTL: السبت on the right.
\newcommand{\dayheaders}{السبت/الأحد/الاثنين/الثلاثاء/الأربعاء/الخميس/الجمعة}

% \daycell{x-index}{y-index}{day-number}{note}
% x-index 0..6 (0 = rightmost = السبت in RTL), y-index 0..5 (0 = top row)
\newcommand{\daycell}[4]{%
  \node[draw=calaccent!60,line width=0.4pt,fill=white,
        minimum width=\cellw cm,minimum height=\cellh cm,
        anchor=north east,
        inner sep=0pt,outer sep=0pt]
    at ($(calgrid.north east)+(-#1*\cellw cm,-#2*\cellh cm)$) (cell-#1-#2) {};
  % day number (top-right corner inside cell, RTL)
  \node[anchor=north east,inner sep=4pt,font=\large\bfseries\color{calhead}]
    at (cell-#1-#2.north east) {#3};
  % notes area
  \node[anchor=north,inner sep=4pt,text width={\cellw cm-0.6cm},
        align=right,font=\footnotesize]
    at ($(cell-#1-#2.north)+(0,-0.7cm)$) {#4};
}

% Empty cell (no day number) for grid padding
\newcommand{\emptycell}[2]{%
  \node[draw=calaccent!30,line width=0.4pt,fill=callight,
        minimum width=\cellw cm,minimum height=\cellh cm,
        anchor=north east,inner sep=0pt,outer sep=0pt]
    at ($(calgrid.north east)+(-#1*\cellw cm,-#2*\cellh cm)$) {};
}

\begin{document}

% ============================================================
%  Month: شهر تجريبي (Sample Month)
%  TODO: غيّر اسم الشهر والسنة، وعدّل أرقام الأيام والملاحظات.
%  الترتيب: العمود 0 = السبت (يمين) ... العمود 6 = الجمعة (يسار).
%  ابدأ اليوم الأول في الخانة المناسبة حسب أول يوم في الشهر.
% ============================================================

\begin{center}
{\fontsize{26pt}{32pt}\selectfont\bfseries\color{calaccent}
شهر تجريبي ١٤٤٧هـ\par}
\vspace{0.05cm}
{\normalsize\color{calhead} التقويم الشهري — الاسم\par}
\vspace{0.2cm}
\begin{tikzpicture}
  \draw[calaccent,line width=1pt] (0,0) -- (\gridw cm,0);
\end{tikzpicture}
\end{center}

\vspace{0.1cm}

\begin{center}
\begin{tikzpicture}
  % Reference origin at top-right of grid (RTL)
  \coordinate (calgrid) at (0,0);

  % ---- Header row (day names) ----
  \foreach \i/\name in {0/السبت,1/الأحد,2/الاثنين,3/الثلاثاء,4/الأربعاء,5/الخميس,6/الجمعة} {
    \node[fill=calhead,text=white,minimum width=\cellw cm,minimum height=0.9cm,
          anchor=north east,inner sep=0pt,font=\large\bfseries]
      at ($(calgrid.north east)+(-\i*\cellw cm,0)$) {\name};
  }

  % Shift grid origin below header
  \coordinate (calgrid) at ($(calgrid)+(0,-0.9cm)$);

  % ---- Day cells (sample: month starts on Sunday => col 1) ----
  % Row 0
  \emptycell{0}{0}
  \daycell{1}{0}{1}{}
  \daycell{2}{0}{2}{}
  \daycell{3}{0}{3}{اجتماع}
  \daycell{4}{0}{4}{}
  \daycell{5}{0}{5}{}
  \daycell{6}{0}{6}{}
  % Row 1
  \daycell{0}{1}{7}{}
  \daycell{1}{1}{8}{موعد طبي}
  \daycell{2}{1}{9}{}
  \daycell{3}{1}{10}{}
  \daycell{4}{1}{11}{}
  \daycell{5}{1}{12}{}
  \daycell{6}{1}{13}{}
  % Row 2
  \daycell{0}{2}{14}{}
  \daycell{1}{2}{15}{}
  \daycell{2}{2}{16}{}
  \daycell{3}{2}{17}{}
  \daycell{4}{2}{18}{}
  \daycell{5}{2}{19}{}
  \daycell{6}{2}{20}{}
  % Row 3
  \daycell{0}{3}{21}{}
  \daycell{1}{3}{22}{}
  \daycell{2}{3}{23}{}
  \daycell{3}{3}{24}{}
  \daycell{4}{3}{25}{}
  \daycell{5}{3}{26}{}
  \daycell{6}{3}{27}{}
  % Row 4
  \daycell{0}{4}{28}{}
  \daycell{1}{4}{29}{}
  \daycell{2}{4}{30}{}
  \emptycell{3}{4}
  \emptycell{4}{4}
  \emptycell{5}{4}
  \emptycell{6}{4}
  % Row 5 (spare, empty)
  \emptycell{0}{5}
  \emptycell{1}{5}
  \emptycell{2}{5}
  \emptycell{3}{5}
  \emptycell{4}{5}
  \emptycell{5}{5}
  \emptycell{6}{5}

\end{tikzpicture}
\end{center}

\end{document}
